\begin{table}[htbp]
\centering
\caption{Heterogeneity by arXiv Category}
\label{tab:heterogeneity}
\begin{adjustbox}{max width=\textwidth}
\begin{tabular}{lccc}
\hline\hline
& (1) & (2) & (3) \\
& Pooled & stat.ML & cs.CV \\
\hline
\multicolumn{4}{l}{\textit{Panel A: Log(3yr Citations + 1)}} \\
After cutoff & $-1.086$ & $-0.929$ & $-1.243$ \\
& (0.733) & (1.568) & (2.156) \\
& $[0.139]$ & $[0.553]$ & $[0.564]$ \\
Eff.\ $N$ & 86 & 17 & 22 \\
\hline
\multicolumn{4}{l}{\textit{Panel B: Any Frontier Citation}} \\
After cutoff & $-0.0183$ & $-0.0214$ & $-0.0128$ \\
& (0.0524) & (0.0643) & (0.0478) \\
& $[0.726]$ & $[0.739]$ & $[0.789]$ \\
Eff.\ $N$ & 86 & 17 & 22 \\
\hline\hline
\end{tabular}
\end{adjustbox}
\begin{minipage}{0.95\textwidth}
\footnotesize
\textit{Notes:} Column (1) reports the pooled estimate across all categories (identical to Tables 2--3). Columns (2)--(3) report category-specific estimates for stat.ML and cs.CV, the two categories with sufficient effective sample sizes ($N \geq 15$) for separate estimation. Categories cs.AI, cs.CL, and cs.LG are excluded from the category-specific analysis: cs.AI dominates the pooled sample (over 80\% of observations in the RDD bandwidth are classified as cs.AI, making the cs.AI-specific estimate essentially identical to the pooled estimate), while cs.CL and cs.LG have extremely small effective samples ($N < 15$) that produce unreliable estimates. Robust bias-corrected standard errors in parentheses; $p$-values in brackets.
\end{minipage}
\end{table}
